\chapter{Grundlagen}
\label{ch:grundlagen}

In diesem Kapitel werden die für das Verständnis der Arbeit notwendigen
Grundlagen beschrieben. Zunächst wird die verwendete Hardware vorgestellt,
anschließend die zugrunde liegende Software-Architektur des MXCK-Systems
und zuletzt der eingesetzte Follow-the-Gap-Algorithmus erläutert.

% ============================================================
\section{Hardware: Das MXCarkit}
% ============================================================

Das MXCarkit ist ein miniaturisiertes autonomes Versuchsfahrzeug,
das im Rahmen von Lehrveranstaltungen an der Technischen Hochschule
Mittelhessen eingesetzt wird. Es dient als Plattform für die Erprobung
autonomer Fahrfunktionen in realen Umgebungen.

\begin{table}[H]
\centering
\caption{Technische Spezifikationen des MXCarkit}
\label{tab:specs}
\begin{tabular}{ll}
\toprule
\textbf{Eigenschaft} & \textbf{Wert} \\
\midrule
Fahrzeugbreite     & $> 30$\,cm \\
Fahrzeuglänge      & $> 50$\,cm \\
Antrieb            & Bürstenloser Elektromotor \\
Motorcontroller    & VESC (Vedder Electronic Speed Controller) \\
Lenkung            & Ackermann-Geometrie \\
Recheneinheit      & NVIDIA Jetson Nano \\
LiDAR-Sensor       & RPLIDAR A1/A2 \\
Betriebssystem     & Ubuntu 18.04 (L4T) \\
Middleware         & ROS~2 (Foxy / Humble) \\
\bottomrule
\end{tabular}
\end{table}

\bildplatzhalter{Foto des MXCarkit von vorne und seitlich mit beschrifteten
Komponenten (LiDAR, Jetson Nano, VESC, Räder)}{Das MXCarkit -- Versuchsfahrzeug
mit RPLIDAR A1/A2 und Jetson Nano}

\subsection{NVIDIA Jetson Nano}

Der Jetson Nano ist ein eingebetteter Einplatinencomputer der Firma NVIDIA
mit einer quad-core ARM Cortex-A57 CPU und einer integrierten 128-core
Maxwell-GPU. Im Kontext dieser Arbeit übernimmt er die vollständige
Verarbeitungskette von der Sensordatenaufnahme über die ROS~2-Middleware
bis hin zur Ausgabe von Steuerkommandos an den VESC.

\subsection{LiDAR-Sensor: RPLIDAR A1/A2}

Der RPLIDAR A1/A2 ist ein kostengünstiger 2D-LiDAR-Sensor, der per
Triangulationsverfahren Distanzmessungen in einem vollständigen
$360°$-Sichtfeld liefert. Relevante Parameter für diese Arbeit:

\begin{table}[H]
\centering
\caption{Technische Kenndaten des RPLIDAR A1/A2}
\label{tab:lidar}
\begin{tabular}{ll}
\toprule
\textbf{Parameter} & \textbf{Wert} \\
\midrule
Scanrate              & $\approx 11$\,Hz \\
Messpunkte/Umdrehung  & $\approx 1800$ \\
Winkelauflösung       & $\approx 0{,}35°$ \\
Minimalmessreichweite & $0{,}15$\,m \\
Maximalmessreichweite & $12$--$16$\,m \\
Sichtfeld             & $360°$ \\
Schnittstelle         & USB (seriell) \\
\bottomrule
\end{tabular}
\end{table}

Der Sensor ist auf dem Fahrzeugdach montiert und publiziert Messdaten
über den ROS~2-Topic \texttt{/scan} im Format
\texttt{sensor\_msgs/msg/LaserScan}. Der Scan beginnt bei
\texttt{angle\_min}\,$= -\pi$ und endet bei \texttt{angle\_max}\,$= +\pi$.

\textbf{Wichtige Konvention:} Im Koordinatensystem des RPLIDAR entspricht
der Winkel $0°$ \emph{nicht} automatisch der Fahrzeugvorderseite.
Die tatsächliche Zuordnung hängt von der physischen Montageposition
des Sensors ab und muss kalibriert werden (siehe Abschnitt~\ref{sec:kalibrierung}).

\subsection{Ackermann-Steuermodell}

Das Fahrzeug verwendet die Ackermann-Lenkgeometrie, bei der Vorder- und
Hinterräder unterschiedliche Lenkwinkel aufweisen, um schlupffreies
Kurvenfahren zu ermöglichen. Steuerkommandos werden im ROS~2-Format
\texttt{ackermann\_msgs/msg/AckermannDriveStamped} mit den Feldern
\texttt{steering\_angle} (Lenkwinkel in Radiant) und
\texttt{speed} (Geschwindigkeit in m/s) übergeben.

% ============================================================
\section{Software-Architektur: MXCK2-Stack}
% ============================================================

\subsection{ROS~2 und Docker-Container}

Das MXCarkit verwendet ROS~2 (Robot Operating System 2) als Middleware.
Die gesamte Softwareumgebung ist in Docker-Containern gekapselt, was
eine reproduzierbare Deploymentumgebung gewährleistet. Für diese Arbeit
relevante Container sind:

\begin{table}[H]
\centering
\caption{Docker-Container-Übersicht auf dem Jetson Nano}
\label{tab:containers}
\begin{tabularx}{\textwidth}{llX}
\toprule
\textbf{Container} & \textbf{ROS~2-Version} & \textbf{Inhalt} \\
\midrule
\texttt{mxck2\_control}     & Humble & Fahrzeugsteuerung, VESC-Treiber, LiDAR-Treiber \\
\texttt{mxck2\_development} & Foxy   & Entwicklungsumgebung, FTG-Stack (diese Arbeit) \\
\texttt{mxck2\_lidar}       & Humble & Alternativer LiDAR-Treiber \\
\bottomrule
\end{tabularx}
\end{table}

Der Workspace \texttt{/mxck2\_ws} im Entwicklungscontainer ist über ein
Docker-Volume direkt aus dem Host-Verzeichnis
\texttt{/home/mxck/development\_ws} eingehängt, sodass Codeänderungen
ohne Neustart des Containers wirksam werden.

\subsection{Basisstack: mxck2\_ws}

Der vorhandene Basisstack (Repository: \texttt{william-mx/mxck2\_ws})
stellt die grundlegende Fahrzeugfunktionalität bereit: manuelle
RC-Steuerung, URDF-Fahrzeugbeschreibung, Kamerabilder und Basistreiber.
Er definiert außerdem den TF-Baum (Transform Tree) des Fahrzeugs mit
den Koordinatenrahmen \texttt{base\_link} (Fahrzeugmittelpunkt) und
\texttt{laser} (LiDAR-Sensor).

\bildplatzhalter{Screenshot des ROS~2-Topic-Graphs (z.B. aus Foxglove)
mit allen aktiven Knoten und Verbindungen}{ROS~2-Knoten- und Topic-Graph
des vollständigen FTG-Stacks}

% ============================================================
\section{Der Follow-the-Gap-Algorithmus}
% ============================================================

\subsection{Grundprinzip}

Der Follow-the-Gap-Algorithmus (FTG) ist ein reaktives Planungsverfahren
für autonome Fahrsysteme, das ohne globale Karte oder Pfadplanung auskommt.
Er operiert ausschließlich auf aktuellen Sensordaten und berechnet in
Echtzeit eine kollisionsfreie Fahrtrichtung.

Das Grundprinzip lässt sich in drei Schritten zusammenfassen:
\begin{enumerate}
    \item Identifiziere alle Hindernisse im Sichtfeld des Sensors.
    \item Finde die größte \emph{Lücke} (freien Winkelbereich) zwischen benachbarten Hindernissen.
    \item Steuere das Fahrzeug in Richtung des Mittelpunkts dieser Lücke.
\end{enumerate}

\begin{equation}
    \theta_{\text{Ziel}} = \frac{\theta_{\text{links}} + \theta_{\text{rechts}}}{2}
    \label{eq:gap_center}
\end{equation}

Dabei sind $\theta_{\text{links}}$ und $\theta_{\text{rechts}}$ die
Winkel der linken und rechten Begrenzung der gewählten Lücke.

\subsection{Klassisches F1TENTH-FTG vs.\ CTU-FTG}

Im Rahmen dieser Arbeit wurde eine modifizierte Variante des FTG-Algorithmus
verwendet, die ursprünglich von der Tschechischen Technischen Universität Prag
(CTU) entwickelt wurde. Im Unterschied zur klassischen F1TENTH-Implementierung
arbeitet diese Variante nicht direkt auf Laserscan-Rohdaten, sondern auf
einer vorverarbeiteten Liste kreisförmiger Hindernisse.

\begin{table}[H]
\centering
\caption{Vergleich: Klassisches F1TENTH-FTG vs. CTU-FTG}
\label{tab:ftg_compare}
\begin{tabularx}{\textwidth}{lXX}
\toprule
\textbf{Merkmal} & \textbf{Klassisches F1TENTH-FTG} & \textbf{CTU-FTG (diese Arbeit)} \\
\midrule
Eingabe          & Rohlaserscan (\texttt{LaserScan})   & Kreishindernisliste (\texttt{ObstaclesStamped}) \\
Sicherheitsradius & Explizit um jeden Messpunkt ($\approx$\,halbe Fahrzeugbreite) & Implizit durch Hindernisradius (0{,}01\,m) \\
Lückenbewertung  & Tiefster freier Strahl             & Größte Winkelbreite \\
Implementierung  & Python oder C++                    & C++ \\
Geschwindigkeit  & Oft direkt integriert              & Separater Adapterknoten \\
\bottomrule
\end{tabularx}
\end{table}

\subsection{Algorithmus-Schritte im Detail}

Der CTU-FTG-Kern (\texttt{follow\_the\_gap\_v0}) durchläuft pro
Planungszyklus folgende Schritte:

\begin{enumerate}[leftmargin=*, label=\textbf{Schritt \arabic*.}]
    \item \textbf{Eingabe konvertieren:} Kreisförmige Hindernisse aus
    \texttt{ObstaclesStamped} in interne \texttt{Obstacle}-Objekte umwandeln.
    \item \textbf{Filtern:} Hindernisse außerhalb des Sichtfelds oder
    jenseits der maximalen Relevanzreichweite verwerfen.
    \item \textbf{Lücken konstruieren:} Zwischen je zwei benachbarten,
    nicht-überlappenden Hindernissen ein \texttt{Gap}-Objekt anlegen.
    \item \textbf{Beste Lücke wählen:} Die Lücke mit der größten
    Winkelbreite auswählen, sofern sie breit genug für das Fahrzeug ist.
    \item \textbf{Zielwinkel berechnen:} Mittelpunktwinkel der Lücke
    nach Gleichung~\eqref{eq:gap_center} berechnen.
    \item \textbf{Fallback:} Falls keine Lücke gefunden wird, Corner-Following
    aktivieren und \texttt{/gap\_found\,=\,false} publizieren.
\end{enumerate}

\subsection{Hindernisrepräsentation: obstacle\_substitution}

Ein wesentlicher Bestandteil des Systems ist der Knoten
\texttt{obstacle\_substitution\_node}, der jeden einzelnen
LiDAR-Messstahl in ein kreisförmiges Hindernis mit Radius~$r = 0{,}01$\,m
umwandelt. Die kartesischen Koordinaten werden aus den Polarkoordinaten
des Scans berechnet:

\begin{equation}
    x = d \cdot \cos(\theta), \quad y = d \cdot \sin(\theta)
    \label{eq:polar2cart}
\end{equation}

wobei $d$ die gemessene Distanz und $\theta$ der zugehörige Scanwinkel ist.
